\subsection{A uniform genus-two criterion for the prescribed point}

\begin{claim}{B.141}{Proved}{Canonical-point classes and the unique potentially bad twist}
Let $C/k$ be a smooth proper geometrically connected curve of genus
two over $k=\mathbf F_q$ of characteristic $p$, let $P\in C(k)$, and put
\[
 J=J_C,\qquad \delta=[K_C-2P]\in J(k).
\]
Let $\ell\ne p$ be an odd prime for which
$f=\operatorname{ord}_{\ell}(p)$ is even.
For every nonzero additive character
$\chi:J(k)\twoheadrightarrow\mathbf Z/\ell$ with
$\chi(\delta)\ne0$, the associated constant cyclic
$P$-split Lang cover $D_\chi\to C$ has
\[
                  p\nmid \frac{\#J_{D_\chi}(k)}{\#J(k)}.
                                                            \tag{B.141.1}
\]
In particular, if $\delta\notin\ell J(k)$, this prescribed-point
class-number witness exists already over $k$.
If $r=\dim_{\mathbf F_\ell}J(k)/\ell J(k)$, this condition gives
at least $\ell^{r-1}$ distinct pointed covers of degree $\ell$,
where an identification of the deck group with
$\mathbf Z/\ell$ is not retained.

More precisely, fix any nonzero $\chi$, put $a=\chi(\delta)$,
and index its constant twists by $t\in\mathbf Z/\ell$ so that
a closed point of degree $m$ acquires the added Frobenius
translation $mt$.  Among these $\ell$ twists, the only
possibly bad twist is
\[
                              2t+a=0.                       \tag{B.141.2}
\]
This twist need not be bad.  In particular, the number of bad
twists is at most one, regardless of the index of
$\langle p\rangle$ in $(\mathbf Z/\ell)^\times$.

If $\chi(\delta)\ne0$, the same cover after base change to
every extension $k_n/k$ of degree $n$ prime to $\ell$ is
again a $P$-split class-number witness.
The sufficient condition here is the nonvanishing of the class
of $\delta$ in $J(k)/\ell J(k)$ at the stated field; nonzero
$\ell$-primary torsion by itself is not this condition.

A nonempty explicit family is
\[
 C_p:\ y^2=x^6+1,\qquad P=(0,1),\qquad
          p\ge5,\quad p\equiv2\text{ or }5\pmod9 .
                                                            \tag{B.141.10}
\]
Here one may take $\ell=3$.
For $p=5$, the character specified in the proof gives a cover
of genus four with
\[
  \#J_{C_5}(\mathbf F_5)=36,\qquad
  \#J_D(\mathbf F_5)=432,\qquad
  \frac{\#J_D(\mathbf F_5)}{\#J_{C_5}(\mathbf F_5)}=12.
                                                            \tag{B.141.11}
\]
\end{claim}
\begin{proof}
We give the pointed construction and the complete character calculation.
The Lang map $F_q-1:J\to J$ is a separable isogeny with
constant kernel $J(k)$.
Factoring its source by $\ker\chi$ gives a separable isogeny
with constant cyclic kernel of order $\ell$.
Its pullback under $j_P:R\mapsto[R-P]$ has a chosen rational
point over $P$, and all of that fibre splits.
The pullback is geometrically connected: the Abel--Jacobi map
induces the isomorphism on prime-to-$p$ geometric abelian
fundamental groups, so the nonzero cyclic quotient remains
nonzero on the curve.  Equivalently, pullback along $j_P$
identifies the degree-zero $\ell$-torsion of the dual Jacobian
with that of $J$, by its principal polarization.
This is the same Lang and Kummer construction used in B.2;
no point-section separation theorem is used.

Choose a primitive $\ell$-th root of unity $\zeta$ in characteristic
zero.  For $i\in(\mathbf Z/\ell)^\times$ write
\[
       S_i=\sum_{R\in C(k)}\zeta^{\,i\chi([R-P])},
                     \qquad a=\chi(\delta).
\]
Expanding the unramified character Euler product by effective
divisors gives
\[
                        L_i(T)=1+S_iT+q\zeta^{ia}T^2.
                                                            \tag{B.141.3}
\]
Here is a direct Riemann--Roch verification, including the last
coefficient.  Tensoring by $\mathcal O(-dP)$ identifies
$\operatorname{Pic}^d(C)(k)$ with $J(k)$.
Every degree-two class other than the canonical class has one
effective divisor over $k$, while the canonical class has
$q+1$ of them: the dimensions of sections are respectively
one and two.  Summing the nontrivial character over all of
$J(k)$ gives zero, leaving the coefficient
$q\zeta^{i\chi([K_C-2P])}$.
For $d\ge3$, Riemann--Roch gives the same section dimension
$d-1$ in every class, so the number
$(q^{d-1}-1)/(q-1)$ of effective divisors in each class is
constant and the character sum is zero.
The degree-one coefficient is the stated sum over $C(k)$.
The rational point $P$ ensures these Picard classes are represented
by line bundles over $k$, so no descent obstruction enters this count.

For the twist indexed by $t$, the corresponding factor is
\[
              L_{i,t}(T)=
       1+\zeta^{it}S_iT+q\zeta^{i(a+2t)}T^2.
                                                            \tag{B.141.4}
\]
Indeed, on a closed point of degree $m$, the character value
is multiplied by $\zeta^{imt}$.
The regular-representation factorization of each closed-point
Euler factor gives
\[
 \frac{Z(D_t,T)}{Z(C,T)}
       =\prod_{i=1}^{\ell-1}L_{i,t}(T),\qquad
 \frac{\#J_{D_t}(k)}{\#J(k)}
       =\prod_{i=1}^{\ell-1}L_{i,t}(1).
                                                            \tag{B.141.5}
\]
The second product is an integer, since it is a rational algebraic
integer, and it equals the positive displayed ratio.
Thus the class-number criterion used here is precisely the
nonvanishing of every factor modulo $p$.

The genus-two canonical pencil is the hyperelliptic pencil,
with an involution $\iota$ defined over $k$ satisfying
$R+\iota R\sim K_C$.
For completeness, the pencil is base-point-free because a
degree-one divisor cannot have two independent sections on a
positive-genus curve; Riemann--Roch then gives
$h^0(K_C-R)=1<h^0(K_C)=2$.
It defines a degree-two map to the projective line.
That map is separable, including in characteristic two:
a purely inseparable degree-two extension of a rational
function field over the perfect field $k$ is rational.
Its nontrivial involution is therefore defined over $k$,
and permutes $C(k)$.
Applying this permutation to the character sum gives the exact identity
\[
                       S_{-i}=\zeta^{-ia}S_i.                \tag{B.141.6}
\]

Reduce $\mathbf Z[\zeta]$ into an algebraic closure of
$\mathbf F_p$.
The image $\bar\zeta$ still has order $\ell$.
Since $f$ is even and $(\mathbf Z/\ell)^\times$ is cyclic,
\[
                             p^{f/2}\equiv-1\pmod\ell.
                                                            \tag{B.141.7}
\]
The coefficients of every $S_i$ are rational integers.
Consequently, if the reduction of $L_{i,t}(1)$ is zero,
then raising that equality to the $p^{f/2}$-th power shows
that the reduction of $L_{-i,t}(1)$ is also zero.
Writing these two equalities out using $q=0$ in the residue field,
\[
        \overline S_i=-\bar\zeta^{-it},
                     \qquad
        \overline S_{-i}=-\bar\zeta^{it}.
\]
Together with (B.141.6), they imply
$\bar\zeta^{\,i(a+2t)}=1$.
Both $i$ and $2$ are invertible modulo $\ell$.
It follows that $a+2t=0$.
This proves (B.141.2) simultaneously for all primitive characters.
The $P$-split twist is $t=0$, so $a\ne0$ excludes every vanishing
factor in (B.141.5), proving (B.141.1).

The dual of the finite vector space $J(k)/\ell J(k)$ separates
its nonzero elements.  If $\delta$ is nonzero in that quotient,
there are $(\ell-1)\ell^{r-1}$ linear functionals nonzero on it.
Each is surjective.  The nonzero scalar multiples have the same
kernel and give the same pointed cover, and different kernels give
different geometric cyclic quotients.
Dividing by $\ell-1$ proves the count of covers.

Finally let $k_n/k$ have degree $n$ prime to $\ell$.
The character of the base-changed Lang cover is
\[
       \chi_n=\chi\circ\operatorname{Nm}_{k_n/k}:J(k_n)
                              \longrightarrow\mathbf Z/\ell.
                                                            \tag{B.141.8}
\]
This follows either from closed-point Artin symbols or from
$F_q^n-1=(F_q-1)(1+F_q+\cdots+F_q^{n-1})$ in the Lang construction.
On a class rational over $k$, the norm is multiplication by $n$.
Thus $\chi_n$ is nonzero, and for the extended canonical-point class,
\[
                            \chi_n(\delta)=n\chi(\delta)\ne0.
\]
The same proof over $k_n$, whose characteristic is still $p$,
gives the claimed persistence for this very cover.
In particular the point stays split by base change.
If instead $\delta\in\ell J(k)$, that divisibility persists under
every field extension.  Therefore this sufficient test cannot
first become applicable after enlarging the constants.

\medskip\noindent\textbf{The explicit whole family and its finite certificate.}
For the primes in (B.141.10), the polynomial $x^6+1$ is squarefree.
Its smooth proper double cover of the projective line has genus
two and two rational points $\infty_+,\infty_-$ above infinity.
The canonical class is their sum.
There is a degree-two morphism
\[
 \pi:C_p\longrightarrow E_p,\qquad
 E_p:\ v^2=u^3+1,\qquad (u,v)=(x^2,y).
                                                            \tag{B.141.12}
\]
Both points at infinity map to the origin $O$ of $E_p$.
The point $R=(0,1)=\pi(P)$ has exact order three:
the tangent line $v=1$ meets the cubic at $R$ with multiplicity
three, so its divisor gives $3R=O$, and $R\ne O$.
Consequently pushforward of degree-zero divisors gives
\[
                         \pi_*(\delta)=-2R=R.
                                                            \tag{B.141.13}
\]
Cubing permutes $\mathbf F_p$ because $p\equiv2\pmod3$.
The usual quadratic-character count, followed by that permutation,
therefore gives
\[
 \#E_p(\mathbf F_p)
   =p+1+\sum_{u\in\mathbf F_p}
                  \left(\frac{u^3+1}{p}\right)
   =p+1 .
\]
The assumptions modulo nine give $v_3(p+1)=1$.
In a finite abelian group whose order has three-adic valuation one,
a nonzero point of order three is not divisible by three.
Thus $R\notin3E_p(\mathbf F_p)$.
If $\delta$ belonged to $3J_{C_p}(\mathbf F_p)$, applying the norm
$\pi_*$ would contradict (B.141.13).
This proves $\delta\notin3J_{C_p}(\mathbf F_p)$ throughout the
family.  Since $\operatorname{ord}_3(p)=2$, the criterion already
proved applies to each member, and its persistence applies over
every extension of degree prime to three.
This family conclusion uses the displayed algebraic argument,
not a finite search.

At $p=5$, let $T=(4,0)\in E_5(\mathbf F_5)$.
The six-element elliptic group is generated by $R$ of order
three and $T$ of order two.
Choose $\chi_E(R)=1$, $\chi_E(T)=0$, and put
$\chi=\chi_E\circ\pi_*:J_{C_5}(\mathbf F_5)\to\mathbf Z/3$.
Equation (B.141.13) gives $\chi(\delta)=1$, so this character is
surjective.  Its values on the pointed curve classes are
\[
\begin{array}{c|rrrrrr}
 Q&(0,1)&(0,4)&(2,0)&(3,0)&\infty_+&\infty_-\\ \hline
 \chi([Q-P])&0&1&2&2&2&2 .
\end{array}                                                  \tag{B.141.14}
\]
For a primitive cube root $\zeta$, the two nontrivial character
factors are therefore
\[
  L_1(T)=1+3\zeta^2T+5\zeta T^2,\qquad
  L_2(T)=1+3\zeta T+5\zeta^2T^2.
\]
Their product is
\[
             1-3T+4T^2-15T^3+25T^4 ,
                                                            \tag{B.141.15}
\]
whose value at one is twelve.
Direct enumeration gives
$\#C_5(\mathbf F_5)=6$ and $\#C_5(\mathbf F_{25})=46$.
There is also a short independent check of the latter count.
The sixth powers in $\mathbf F_{25}^{\times}$ form
$\boldsymbol\mu_4=\mathbf F_5^\times$, with six preimages
per value.  The values of $1+\boldsymbol\mu_4$ are
$2,3,4,0$; every nonzero element of $\mathbf F_5$ is a
square in $\mathbf F_{25}$.
Thus the affine count is $2+6(2+2+2+1)=44$, and the two
points at infinity give forty-six.
The first two Newton identities and the reciprocal genus-two
Weil polynomial give
$\#J_{C_5}(\mathbf F_5)=(6^2+46)/2-5=36$.
Together with (B.141.5), this proves (B.141.11).

The complete finite certificate is
\path{finite_checks/code/b141_canonical_example.py}, with its
retained JSON and output.
It enumerates both curve point sets, uses
$\mathbf F_{25}=\mathbf F_5[u]/(u^2-2)$, retains and checks the
whole elliptic addition and character tables, verifies
$R\notin3E_5(\mathbf F_5)$ and all six pointed character values,
and computes (B.141.15) both by quadratic-integer multiplication
and by an integer resultant.
The cover itself is the explicit Lang quotient based at $P$
for the fully specified character $\chi$.
\end{proof}

\paragraph{Exact reduction of the remaining B1 locus.}
This proves B1 for every genus-two marked curve satisfying
$[K_C-2P]\notin\ell J_C(k)$ for at least one odd prime
$\ell\ne p$ with even $\operatorname{ord}_{\ell}(p)$.
It is a condition on the entire marked curve and its given field,
and includes every prime-to-$\ell$ constant extension of each
certified pointed cover.
A counterexample to B1 in genus two must therefore satisfy
\[
 [K_C-2P]\in\ell J_C(k)
 \quad\text{for every odd }\ell\ne p
          \text{ with even }\operatorname{ord}_{\ell}(p).
                                                            \tag{B.141.9}
\]
Only primes dividing $\#J_C(k)$ need be checked in this necessary
condition.  It leaves every Weierstrass marking, for which the
canonical-point class is zero, as well as other possible markings.
It proves neither the existence of a witness nor its impossibility
on that remaining locus, and does not settle the full threshold
$p>(2g-2)^2$.

The preceding B.28 bounds the number of bad twists by
$(\ell-1)/f$ in genus two; the bound here is one whenever $f$ is
even, and identifies the only potentially bad twist relative to
the prescribed point.  B.23 has canonical-point class zero at its
displayed marking and does not supply this sufficient family.
The explicit characteristic-two examples in B.21--B.27 and A.50
do not contain the family (B.141.10).
The proof above is a direct character calculation; no priority
or claim of a complete literature classification is made.

\paragraph{Check record.}
Written and checked by GPT-6 Astra Pro, finite-checks reader,
18 November 2026; the separate effort setting is not exposed.
The complete B.23 and B.28--B.29 texts were reread to check the
scope of the predecessor statements.
The adjacent B.21--B.22 statements were also reread when checking
that the explicit family is outside those finite examples.
The root and configuration readers independently passed the
complete criterion proof before the explicit-family addition;
the configuration reader also passed the entire addition and
the finite certificate, and supplied the independent
$\mathbf F_{25}$ count above.
These are recorded as separate check events.
This statement was prepared after checkpoint04 and was not
included in its frozen PDF.
